% Guion de 15 minutos para la presentación v3 de PULSO.
% Compilar desde docs/latex/guion_v3_15min:
%   latexmk -lualatex guion_15min.tex
\documentclass[11pt,a4paper,oneside]{article}

\usepackage{fontspec}
\usepackage[spanish,es-noshorthands,es-tabla,es-nodecimaldot]{babel}
\usepackage{amsmath,amssymb}
\setmainfont{TeX Gyre Pagella}
\setsansfont{TeX Gyre Heros}
\setmonofont{DejaVu Sans Mono}[Scale=0.82]
\usepackage[margin=2.3cm,headheight=14pt]{geometry}
\usepackage{booktabs,longtable,array}
\usepackage{xcolor}
\usepackage[most]{tcolorbox}
\usepackage{enumitem}
\usepackage{needspace}
\usepackage{titlesec}
\usepackage{fancyhdr}
\usepackage[hidelinks]{hyperref}

\definecolor{azul}{HTML}{2563A8}
\definecolor{rojo}{HTML}{C2362B}
\definecolor{verde}{HTML}{2F7D5B}
\definecolor{gris}{HTML}{5F6B77}

\titleformat{\section}{\sffamily\bfseries\Large\color{azul}}{}{0em}{}
\titleformat{\subsection}{\sffamily\bfseries\large}{}{0em}{}

\pagestyle{fancy}
\fancyhf{}
\fancyhead[L]{\small\sffamily\color{gris}PULSO · guion de 15 minutos (v3)}
\fancyhead[R]{\small\sffamily\color{gris}\nouppercase{\leftmark}}
\fancyfoot[C]{\small\thepage}
\renewcommand{\headrulewidth}{0.3pt}

\setlength{\parskip}{0.5em}
\setlength{\parindent}{0pt}
\setlist{itemsep=0.15em,topsep=0.25em}
\setlist[description]{font=\normalfont\bfseries,leftmargin=1.2em,labelindent=0em}

% \diapo{número}{título}{duración}{reloj al terminar}{mensaje clave}
\newcommand{\diapo}[5]{%
  \par\needspace{9\baselineskip}\vspace{0.7em}%
  \begin{tcolorbox}[colback=azul!8,colframe=azul,boxrule=0.6pt,arc=2pt,
    left=6pt,right=6pt,top=3pt,bottom=3pt,fontupper=\sffamily]
  \textbf{#1}\quad #2\hfill{\small\color{gris}#3 · termina en \textbf{#4}}\\[2pt]
  {\small\color{azul!70!black}\textbf{Mensaje:} #5}
  \end{tcolorbox}\nopagebreak}
\newtcolorbox{decir}{colback=white,colframe=gris!40,boxrule=0.4pt,arc=2pt,
  left=7pt,right=7pt,top=4pt,bottom=4pt,breakable,
  before upper={\setlength{\parskip}{0.6em}}}

\hypersetup{pdftitle={PULSO: guion de 15 minutos},pdfauthor={Juan Diego Martínez Ayala}}

\begin{document}

\begin{center}
{\sffamily\bfseries\huge\color{azul} PULSO · guion de 15 minutos\par}
\vspace{0.2cm}
{\sffamily\color{gris} Presentación v3 (\texttt{pulso\_presentacion\_v3.pdf}) ·
Juan Diego Martínez Ayala · noviembre de 2026\par}
\end{center}

\section{Cómo usar este guion}

Es la versión corta del guion v3. Usa las mismas 23 diapositivas, pero dice solo lo
esencial de cada una. El texto está pensado para leerse o parafrasearse en voz alta.

\begin{itemize}
\item \textbf{Duración}: unos 13 minutos a 125 palabras por minuto, con 3 segundos por
cambio de diapositiva. Quedan casi 2 minutos de margen.
\item \textbf{Prioridad}: el tiempo se concentra en la idea (diapositiva 2), el modelo
(5 y 6), cómo se mide desde las transacciones (12), el diseño del experimento (14) y el
resultado (15 a 17).
\item \textbf{Pasos rápidos}: en las diapositivas 8, 9, 10, 18, 21 y 23 basta con una o
dos frases. Si se atrasa, en ellas diga solo el mensaje de la barra azul.
\item \textbf{Preguntas}: las respuestas preparadas, la definición de cada variable y las
diapositivas de respaldo están en el guion completo, \texttt{pulso\_guion\_v3.pdf}.
\end{itemize}

\subsection{Plan de tiempos}
\begin{center}
\small
\begin{longtable}{lp{0.64\textwidth}cc}
\toprule
\textbf{N.º} & \textbf{Diapositiva} & \textbf{Dura} & \textbf{Termina} \\
\midrule
\endhead
Portada & PULSO & 0:25 & 0:25 \\
1 & Los pagos forman una red y los problemas viajan por ella & 0:35 & 1:00 \\
2 & No solo importa cuánto se recibe, también cuándo & 0:50 & 1:50 \\
3 & Pregunta de investigación & 0:50 & 2:40 \\
4 & Seis tareas, un solo modelo & 0:30 & 3:10 \\
5 & La ecuación del déficit & 0:55 & 4:05 \\
6 & La ecuación del ritmo de pagos & 0:55 & 5:00 \\
7 & La red amortigua los choques solo si $R<1$ & 0:35 & 5:35 \\
8 & Los ciclos de pago se sincronizan como predice la teoría & 0:15 & 5:50 \\
9 & Las cascadas son máximas con conectividad intermedia & 0:20 & 6:10 \\
10 & El aprendizaje automático predice el colapso en milisegundos & 0:20 & 6:30 \\
11 & Una cartera bancaria en tres capas & 0:35 & 7:05 \\
12 & El banco solo ve transacciones; el ritmo se estima desde ellas & 0:45 & 7:50 \\
13 & Tres niveles de información & 0:30 & 8:20 \\
14 & El experimento compara dos mundos idénticos & 0:40 & 9:00 \\
15 & Cada nivel de información mejora la predicción & 0:40 & 9:40 \\
16 & La mejora es real: aparece con red y desaparece en el control & 0:35 & 10:15 \\
17 & La mejora crece con la intensidad del desfase & 0:25 & 10:40 \\
18 & El score clásico sigue siendo la base & 0:20 & 11:00 \\
19 & Uso práctico: la quiebra de un empleador & 0:30 & 11:30 \\
20 & Qué demuestra y qué no & 0:30 & 12:00 \\
21 & Software reproducible y aplicación web & 0:15 & 12:15 \\
22 & Conclusiones & 0:40 & 12:55 \\
23 & Trabajo futuro & 0:15 & 13:10 \\
Cierre & Gracias & 0:05 & 13:15 \\
\bottomrule
\end{longtable}
\end{center}

\subsection{Puntos de control}
\begin{center}
\begin{tabular}{lll}
\toprule
Al terminar & Debería ir en & Si va atrasado \\
\midrule
Diapositiva 6 & 5:00 & En 7 a 10 diga solo el mensaje de cada barra azul \\
Diapositiva 14 & 9:00 & En 18 diga solo el mensaje \\
Diapositiva 20 & 12:00 & Omita la 21 y la 23 \\
\bottomrule
\end{tabular}
\end{center}

% Texto de cada diapositiva de la presentación v3 en versión de 15 minutos.
% Tiempos: conteo de palabras a 125 palabras por minuto más 3 s por cambio de diapositiva.

% ===================================================================
\section{El problema}

\diapo{Portada}{PULSO}{0:25}{0:25}{Quién presenta y de qué trata.}
\begin{decir}
Buenos días. Presento PULSO, el proyecto del Semillero de Modelamiento de Sistemas
Complejos y Física Computacional. La pregunta es si el momento en que llegan los pagos,
y no solo su monto, ayuda a predecir quién deja de pagar un crédito.
\end{decir}

\diapo{1}{Los pagos forman una red y los problemas viajan por ella}{0:35}{1:00}{Un
problema de liquidez no se queda donde empieza.}
\begin{decir}
Cada persona y cada empresa paga con lo que recibe de otras, así que los pagos forman una
red. En el dibujo, A deja de pagar: B y C reciben menos y, si no tienen reservas, tampoco
le pagan a D. El problema llega a D sin que tenga relación con A. Con pagos inmediatos
como Bre-B, esto ocurre más rápido.
\end{decir}

\diapo{2}{No solo importa cuánto se recibe, también cuándo}{0:50}{1:50}{El ingreso
desfasado genera estrés aunque el monto no cambie.}
\begin{decir}
Esta es la idea central. Las dos líneas son la misma persona, con el mismo salario y la
misma cuota. Arriba, el salario llega antes de que venza la cuota y paga a tiempo. Abajo,
el salario llega tres días después: durante esos días debe la cuota y no tiene con qué
pagarla.

El monto no cambió; cambió la fecha. A eso lo llamamos desfase o desincronización. Si se
repite cada mes y no hay ahorro, termina en mora. Lo que planteamos es que este efecto se
puede medir y sirve para predecir.
\end{decir}

\diapo{3}{Pregunta de investigación}{0:50}{2:40}{Comparamos contra la red estática porque
eso ya existe.}
\begin{decir}
Los bancos predicen el impago con un score que usa datos del cliente: ingreso, deuda e
historial de mora. Nuestra pregunta es si la sincronización de los pagos, medida solo con
las transacciones que el banco ya tiene, mejora esa predicción. Y la comparamos con dos
referencias: el score clásico, y la red estática, es decir, quién le paga a quién y
cuánto, que ya se usa en la literatura. Lo nuevo es medir el ritmo.

La hipótesis se puede refutar: si el desfase no genera déficit, la sincronización no
debe ayudar. Para comprobarlo usamos un mundo de control.
\end{decir}

\diapo{4}{Seis tareas, un solo modelo}{0:30}{3:10}{Cada pieza se valida antes de agregar
la siguiente.}
\begin{decir}
El proyecto se construyó en seis tareas. Las cinco primeras arman el modelo: red y
déficit, redistribución, sincronización, cascadas de impago y predicción con aprendizaje
automático. La sexta lo aplica a un banco. Cada tarea agrega un término, y si se apaga,
el modelo vuelve exactamente al anterior. Así cada pieza se prueba por separado.
\end{decir}

% ===================================================================
\section{El modelo}

\diapo{5}{La ecuación del déficit}{0:55}{4:05}{Cinco efectos en una ecuación.}
\begin{decir}
$x_i$ es el déficit de la entidad $i$: cuánta liquidez le falta. La ecuación dice cómo
cambia, y cada color es un efecto.

Verde: cada entidad corrige su déficit por sí misma, por ejemplo con ahorro; alfa es qué
tan rápido. Rojo: el contagio; si quienes le pagan tienen déficit, le pagan menos.
$W_{ij}$ es cuánto le paga $j$ a $i$ y beta la intensidad. Azul: la liquidez se
redistribuye entre entidades sin crearse ni destruirse. Morado: el desfase, que explico
en la siguiente diapositiva. Y al final, los choques externos.

Si el déficit supera el colchón de la entidad, la entidad deja de pagar: eso es un
impago.
\end{decir}

\diapo{6}{La ecuación del ritmo de pagos}{0:55}{5:00}{El término $\gamma k_i(1-r_i)$ une el
ritmo con la liquidez.}
\begin{decir}
Cada entidad paga en ciclos, como un reloj que da vueltas, y theta es la posición de la
manecilla. Es el modelo de Kuramoto, el modelo clásico de sincronización en física: cada
entidad tiene su ritmo propio, omega, y lo ajusta al de quienes le pagan con una fuerza
$K$.

La fórmula de la derecha es la coherencia, $r_i$: vale uno si el ingreso llega a tiempo y
menos uno si llega en el peor momento. Por eso el término morado de la ecuación anterior,
uno menos $r_i$, es cero cuando todo llega a tiempo y crece con el desfase: convierte el
desfase en déficit.
\end{decir}

% ===================================================================
\section{Lo que muestra el modelo}

\diapo{7}{La red amortigua los choques solo si $R<1$}{0:35}{5:35}{Con $R=0.9$ el déficit se
apaga; con $R=1.1$ crece.}
\begin{decir}
Primer resultado: un solo número decide si la red absorbe un choque. Es $R$, el contagio
dividido por la corrección, como el número reproductivo de una epidemia. En las gráficas,
cada línea es el déficit de una entidad y la negra es el total. Con $R$ de 0.9 el total
baja; con 1.1 crece sin límite. Coincide con la solución exacta.
\end{decir}

\diapo{8}{Los ciclos de pago se sincronizan como predice la teoría}{0:15}{5:50}{El código
reproduce un resultado exacto.}
\begin{decir}
Segundo: la sincronía crece con el acoplamiento $K$, y en la red completa el código
reproduce el valor teórico de Kuramoto, 0.7071 contra 0.7076 simulado.
\end{decir}

\diapo{9}{Las cascadas son máximas con conectividad intermedia}{0:20}{6:10}{Más conexiones
no siempre significa más riesgo.}
\begin{decir}
Tercero: con umbrales aparecen cascadas de impago. Son máximas con conectividad
intermedia: con pocas conexiones el choque no tiene por dónde pasar, y con muchas cada
entidad reparte la pérdida entre varios pagadores.
\end{decir}

\diapo{10}{El aprendizaje automático predice el colapso en milisegundos}{0:20}{6:30}{El
modelo físico genera los datos; el aprendizaje automático los resume.}
\begin{decir}
Cuarto: con 20\,000 simulaciones entrenamos modelos que predicen si la red colapsa. XGBoost
logra un AUC de 0.98 y responde en 0.2 milisegundos, frente a 0.08 segundos de una
simulación.
\end{decir}

% ===================================================================
\section{Aplicación al riesgo de crédito}

\diapo{11}{Una cartera bancaria en tres capas}{0:35}{7:05}{Se simula una cartera para
comprobar que el método encuentra lo que sabemos que está.}
\begin{decir}
Ahora los nodos son clientes de un banco: 2000 personas y 100 empresas, con 55~\% de
informalidad. Se relacionan por tres capas: la nómina, de la empresa al empleado; los
pagos a proveedores, entre empresas; y las transferencias entre personas.

No hay datos públicos de este tipo, así que simulamos la cartera con un mecanismo
conocido y verificamos que el método lo encuentra, como quien calibra un instrumento.
\end{decir}

\diapo{12}{El banco solo ve transacciones; el ritmo se estima desde ellas}{0:45}{7:50}{La
coherencia estimada coincide con la real.}
\begin{decir}
Este es el paso clave. El banco no ve fases ni déficits: solo ve pagos, con fecha,
pagador, receptor y monto. Con esos pagos estimamos el ciclo de cada cliente: entre un
pago y el siguiente completó una vuelta, así que a mitad de camino está a media vuelta.
Con las fases estimadas calculamos su coherencia con quienes le pagan.

La figura compara lo real con lo estimado: la correlación es mayor que 0.8. Usamos seis
meses para observar y predecimos el impago de los doce meses siguientes.
\end{decir}

\diapo{13}{Tres niveles de información}{0:30}{8:20}{El nivel 3 es el aporte del proyecto.}
\begin{decir}
Comparamos tres niveles de variables, cada uno contiene al anterior. Uno, el score
clásico. Dos, el clásico más la red estática: de quién recibe, cuánto y qué tan
concentrado está su ingreso. Tres, todo lo anterior más la dinámica: coherencia, desfase
y regularidad de los pagos. Mismo modelo, mismos clientes y misma validación.
\end{decir}

\diapo{14}{El experimento compara dos mundos idénticos}{0:40}{9:00}{Si la dinámica
mejorara en el mundo nulo, el método estaría engañando.}
\begin{decir}
Para que el resultado sea creíble simulamos dos mundos con los mismos clientes, la misma
red y el mismo azar. En uno, el desfase y el contagio generan déficit. En el otro, el
mundo nulo, los pagos ocurren igual pero no generan déficit.

La dinámica debe ayudar en el primero y no en el segundo. Si ayudara en el nulo, el
método estaría encontrando patrones falsos. Usamos 20 carteras por mundo, unos 40\,000
clientes.
\end{decir}

% ===================================================================
\section{Resultados}

\diapo{15}{Cada nivel de información mejora la predicción}{0:40}{9:40}{La dinámica agrega
información que la red estática no tiene.}
\begin{decir}
Para medir uso el AUC: la probabilidad de que el modelo ponga a quien cae por encima de
quien no cae. 0.5 es azar y 1 es perfecto.

El score clásico logra 0.769. Con la red estática, 0.784. Con la dinámica, 0.795. Con
XGBoost el resultado final es el mismo. Todo se mide en carteras que el modelo no vio al
entrenar. La dinámica suma encima de la red estática: no repite lo que ya se sabía.
\end{decir}

\diapo{16}{La mejora es real: aparece con red y desaparece en el control}{0:35}{10:15}{Con
red, positiva en 5 de 5 particiones; en el nulo, cero.}
\begin{decir}
Cada fila es la mejora de la dinámica sobre la red estática, con su intervalo de
confianza del 95~\%. Con red, en rojo, la mejora es de 0.011 y 0.014, lejos del cero, y
positiva en las cinco particiones. En el mundo nulo, en gris, es cero.

Este es el resultado más importante: con las mismas variables y el mismo método, la
mejora aparece solo cuando el mecanismo existe.
\end{decir}

\diapo{17}{La mejora crece con la intensidad del desfase}{0:25}{10:40}{La señal viene de
la desincronización, no del contagio.}
\begin{decir}
Para saber de dónde viene la señal variamos solo gamma, la fuerza del desfase, y dejamos
el contagio igual. Con gamma cero la mejora desaparece, aunque el contagio siga activo.
Al subir gamma, la mejora crece hasta 0.029. La señal viene de la desincronización, no
del contagio.
\end{decir}

\diapo{18}{El score clásico sigue siendo la base}{0:20}{11:00}{El modelo híbrido
complementa el score; no lo reemplaza.}
\begin{decir}
En la importancia de las variables, las clásicas, en gris, siguen pesando más: deuda,
ingreso y mora. Las de red, en azul y rojo, agregan información. El modelo complementa al
score, no lo reemplaza.
\end{decir}

\diapo{19}{Uso práctico: la quiebra de un empleador}{0:30}{11:30}{Sirve para pruebas de
estrés por concentración en empleadores.}
\begin{decir}
El modelo también permite simular escenarios. Si quiebra el mayor empleador del banco,
con 74 empleados, el 12~\% de ellos cae de inmediato, y a seis meses el 13~\% frente al
7~\% sin quiebra. En el mundo nulo no pasa nada. Para un banco es una prueba de estrés por
concentración en un empleador.
\end{decir}

\diapo{20}{Qué demuestra y qué no}{0:30}{12:00}{El siguiente paso es una cartera real
anonimizada.}
\begin{decir}
Demuestra que, si el mecanismo existe, se detecta con las transacciones del banco, que
aporta más que la red estática y que desaparece sin el mecanismo. No demuestra que exista
en un banco real, porque los datos son simulados y los parámetros ilustrativos. El
siguiente paso es probarlo con una cartera real anonimizada.
\end{decir}

% ===================================================================
\section{Cierre}

\diapo{21}{Software reproducible y aplicación web}{0:15}{12:15}{Todo resultado se puede
regenerar y explorar.}
\begin{decir}
Todo está implementado en Python, con una aplicación web de nueve páginas, 127 pruebas
automáticas y Docker. El experimento completo se repite con un solo comando en unos 20
segundos.
\end{decir}

\diapo{22}{Conclusiones}{0:40}{12:55}{No solo importa cuánto recibe un cliente, también
cuándo.}
\begin{decir}
En conclusión: un solo modelo une contagio, redistribución, sincronización y cascadas, y
cada parte se validó contra resultados exactos. El número $R$ dice si la red amortigua o
amplifica. Y el resultado principal: la sincronización de los pagos, medida solo con
transacciones, mejora la predicción del impago más allá del score y de la red estática, y
no lo hace cuando el mecanismo no existe.

No solo importa cuánto recibe un cliente, también cuándo.
\end{decir}

\diapo{23}{Trabajo futuro}{0:15}{13:10}{Validación con datos reales.}
\begin{decir}
Como trabajo futuro: probarlo con datos reales, incluir las fechas de cuota y estudiar qué
pasa cuando un banco solo ve a sus clientes.
\end{decir}

\diapo{Cierre}{Gracias}{0:05}{13:15}{Abrir preguntas.}
\begin{decir}
Muchas gracias. Quedo atento a sus preguntas.
\end{decir}


\section{Definiciones mínimas}

Para tener a mano durante la exposición. El detalle está en el guion completo.

\begin{description}
\item[Déficit ($x_i$).] Liquidez que le falta a una entidad.
\item[Colchón.] Déficit máximo que una entidad soporta antes de dejar de pagar.
\item[$W_{ij}$.] Cuánto le paga $j$ a $i$ en cada ciclo.
\item[$\alpha$, $\beta$, $\gamma$.] Velocidad de corrección, intensidad del contagio y
fuerza del desfase.
\item[$R=\beta\rho(W)/\alpha$.] Contagio sobre corrección; $\rho(W)$ es el valor propio
dominante de la red. Menor que 1: la red amortigua.
\item[Fase ($\theta_i$) y ritmo ($\omega_i$).] Punto del ciclo de pagos y velocidad
propia del ciclo.
\item[Kuramoto.] Modelo de osciladores que ajustan su ritmo al de sus vecinos.
\item[Coherencia ($r_i$).] De $-1$ a $1$: qué tan a tiempo llega el ingreso respecto al
ciclo propio.
\item[Score clásico.] Modelo de probabilidad de impago con datos del cliente.
\item[Red estática.] Quién le paga al cliente y cuánto, sin mirar el momento.
\item[Red dinámica.] Coherencia, desfase y regularidad de los pagos.
\item[AUC.] Probabilidad de que el modelo ponga a quien cae por encima de quien no cae.
\item[Intervalo de confianza 95~\%.] Rango que contiene la mejora verdadera con 95~\% de
confianza; si no toca el cero, la mejora es significativa.
\item[Mundo nulo.] Mundo idéntico donde los pagos no generan déficit; sirve de control.
\item[XGBoost.] Modelo de árboles de decisión entrenados en secuencia.
\end{description}

\end{document}
